% SPDX-License-Identifier: Apache-2.0
\section{The Operators}
\label{sec:x-ops}

The operators of a value are Rust's. \code{+} and \code{-} are
same-width and wrapping; \code{\&}, \code{|}, \code{\^{}} and
\code{!} are bitwise; \code{<<} and \code{>>} shift by an integer,
not by a value; and the compares, \code{==}, \code{!=}, \code{<},
\code{<=}, \code{>} and \code{>=}, are unsigned and yield a
\code{bool}. The right operand of an arithmetic or bitwise operator,
and of a compare, is anything that converts, so a literal is written
as a literal: \code{a + 1}, \code{a \& 0x3F}, \code{a >= 128}. A
\code{bool} is a condition wherever a \code{Bit} is, in \code{with!},
\code{mux}, a channel's \code{recv\_if} and a register's
\code{set\_if}, and either drives an output or a register of
\code{Bit}; \code{\&} and \code{|} join the two, and the result is a
\code{Bit}. Rust binds \code{\&} tighter than \code{==}, so a compare
inside a logic expression is parenthesised, and \code{\&\&} and
\code{||} cannot be overloaded, so they stay on \code{bool}. What has
no operator is a method: \code{sra}, the arithmetic shift, and
\code{lt\_signed}, the signed compare, in \code{txhdl::funcs}, and
\code{mul::<M>}, \code{concat::<\_, M>}, \code{sext::<M>} and
\code{zext::<M>}, each with the result width stated.
Each operator lowers to the operator of the same name; a compare in
VHDL is a boolean and is made a bit where a bit is needed. One case
of that is worth saying, because the two languages disagree about
what a value of one bit is. A slice of one bit, \code{a.slice::<7,
1>()}, compared with a number is how a program asks whether that bit
is set. In Verilog it is \code{a[7:7] == 1}, a vector against a
number. In VHDL a value of one bit is a \code{std\_logic} rather than
an array of one element, so it is written \code{a(7) = '1'}: a range
there would be an \code{unsigned}, which has no \code{=} against a
bit literal, and the analyser would call the literal ambiguous. The
lowering writes each the way its language wants it.

The unit below drives seven registers from two inputs under
\code{en}, counts the cycles on which the inputs are equal or ordered
and the first is below 128, counts in a second register the cycles on
which either input has its top bit set, and says on two wires whether
they are equal and whether the first is below. The top bit is read
through a name of its own for the first input and inside the
condition for the second, which are the two places the lowering has
to get right. The lowering is simulated against the run's
trace, Section~\ref{sec:x-checked}.

\begin{figure}[htbp]
\centering
\input{ops_timing.tex}
\caption{The operators: the inputs, the sum, the two counts, and the
two compares.}
\label{fig:x-ops}
\end{figure}

\lstinputlisting[caption={\code{ex\_ops.rs}. Run with \code{bazel run //lib/examples:ex\_ops}.},label={lst:x-ops}]{lib/examples/ex_ops.rs}

\lstinputlisting[style=ascii,caption={What \code{ex\_ops} printed when the build ran it: a line per cycle, then the lowering.},label={lst:x-ops-out}]{out_ops.txt}
